% Acronyms and the style of their list.

\makeglossaries
% Each acronym has a short description, shown only in the list of acronyms.
% Keep them consistent with the definition boxes of the text.
\setabbreviationstyle[acronym]{long-short-desc}
\renewcommand*{\glsxtrlongshortdescsort}{\expandonce\glsxtrorgshort}
% List style: short form on the left, long form and description on the right.
\newlength{\AcronymShortWidth}
\setlength{\AcronymShortWidth}{0.12\linewidth}
\newglossarystyle{acronymdesc}{%
  \setglossarystyle{long}%
  \renewenvironment{theglossary}%
    {\begin{longtable}{@{}p{\AcronymShortWidth}p{\dimexpr\linewidth-\AcronymShortWidth-2\tabcolsep\relax}@{}}}%
    {\end{longtable}}%
  \renewcommand{\glossentry}[2]{%
    \glsentryitem{##1}\glstarget{##1}{\textbf{\glsentryshort{##1}}} &
    \textbf{\glsentrylong{##1}}\newline
    \glossentrydesc{##1}\par\vspace{0.6em}\tabularnewline
  }%
  \renewcommand*{\glsgroupskip}{}% uniform spacing, no gap between letter groups
}
\setglossarystyle{acronymdesc}
% Label of the list of acronyms, referenced from the introduction.
\setglossarypreamble[\acronymtype]{\label{chap:acronyms}}

% General fields
\newacronym[description={The field of computer science concerned with building systems that perform tasks usually associated with human intelligence, such as perception, reasoning, learning and decision-making.}]
  {ai}{AI}{Artificial Intelligence}
\newacronym[description={The subfield of artificial intelligence that develops methods to make the decisions of a model understandable to a human observer (see Section~\ref{subsec:explainable_artificial_intelligence}).}]
  {xai}{XAI}{Explainable Artificial Intelligence}
\newacronym[description={The branch of explainability dedicated to reinforcement learning agents, whose decisions are sequential and depend on the interaction with an environment (see Chapter~\ref{chap:explainability_in_reinforcement_learning}).}]
  {xrl}{XRL}{Explainable Reinforcement Learning}
\newacronym[description={The explainability approach that seeks to explain the behavior of a neural network by analyzing its internal mechanisms (see Definition~\ref{def:mechanistic_explainability}).}]
  {me}{ME}{Mechanistic Explainability}

% Deep learning
\newacronym[description={The family of machine learning methods based on neural networks with many successive layers, trained from data by gradient descent (see Chapter~\ref{chap:deep_learning}).}]
  {dl}{DL}{Deep Learning}
\newacronym[description={A parameterized function made of successive layers of artificial neurons, whose parameters are learned from data.}]
  {nn}{NN}{Neural Network}
\newacronym[description={A neural network made of a sequence of fully connected layers.}]
  {mlp}{MLP}{Multilayer Perceptron}
\newacronym[description={A neural network whose layers apply learned convolution filters, suited to data with a spatial structure such as images.}]
  {cnn}{CNN}{Convolutional Neural Network}
\newacronym[description={A recurrent neural network architecture whose gating mechanism allows information to be retained over long sequences.}]
  {lstm}{LSTM}{Long Short-Term Memory}
\newacronym[description={A neural network trained on large text corpora to predict the next token of a sequence, used for language understanding and generation.}]
  {llm}{LLM}{Large Language Model}
\newacronym[description={An optimization algorithm that updates the parameters of a model in the direction opposite to the gradient of the loss, estimated on a random subset of the training data (see Section~\ref{sec:deep_learning_learning}).}]
  {sgd}{SGD}{Stochastic Gradient Descent}
\newacronym[description={A loss function equal to the average of the squared differences between predictions and targets.}]
  {mse}{MSE}{Mean Squared Error}

% Latent space analysis
\newacronym[description={The internal encoding of an input given by an intermediate layer of a trained neural network (see Definition~\ref{def:representation}).}]
  {lr}{LR}{Latent Representation}
\newacronym[description={The set of all latent representations produced by a layer over a dataset (see Definition~\ref{def:latent_space}).}]
  {ls}{LS}{Latent Space}
\newacronym[description={The specific form in which a concept is instantiated within the latent space, for instance a direction or a cluster (see Definition~\ref{def:concept_structure}).}]
  {cs}{CS}{Concept Structure}
\newacronym[description={An autoencoder trained with a sparsity constraint, used to decompose latent representations into a larger set of more interpretable directions (see Section~\ref{sec:latent_space_analysis_superposition}).}]
  {sae}{SAE}{Sparse Autoencoder}

% Reinforcement learning
\newacronym[description={The learning paradigm in which an agent learns a policy by interacting with an environment so as to maximize the rewards it collects (see Chapter~\ref{chap:reinforcement_learning}).}]
  {rl}{RL}{Reinforcement Learning}
\newacronym[longplural={Markov Decision Processes}, description={The mathematical framework modeling sequential decision-making as an agent interacting with an environment through states, actions, rewards and transitions (see Section~\ref{subsec:markov_decision_processes}).}]
  {mdp}{MDP}{Markov Decision Process}
\newacronym[description={A reinforcement learning algorithm that approximates the action-value function with a deep neural network.}]
  {dqn}{DQN}{Deep Q-Network}
\newacronym[description={An actor--critic reinforcement learning algorithm for continuous action spaces, which learns a deterministic policy.}]
  {ddpg}{DDPG}{Deep Deterministic Policy Gradient}
\newacronym[description={A policy-gradient reinforcement learning algorithm that stabilizes training by limiting the size of each policy update through a clipped objective.}]
  {ppo}{PPO}{Proximal Policy Optimization}
\newacronym[description={A variant of Proximal Policy Optimization in which data collection by the workers and policy updates run asynchronously, to scale training.}]
  {appo}{APPO}{Asynchronous Proximal Policy Optimization}
\newacronym[description={The training of a policy by supervised learning to reproduce the actions of another policy on a dataset of its interactions.}]
  {bc}{BC}{Behavioral Cloning}

% Proposed method and evaluation metrics
\newacronym[description={The method proposed in this thesis, which trains several surrogate policies of an agent and clusters their latent representations to detect concept structures (see Chapter~\ref{chap:clustering_agent_representations}).}]
  {car}{CAR}{Clustering Agent Representation}
\newacronym[description={A measure of the agreement between two partitions of the same set of points, computed over all pairs of points (see Section~\ref{subsec:adjusted_rand_index}).}]
  {ri}{RI}{Rand Index}
\newacronym[description={The Rand Index corrected for chance, equal to $0$ for random partitions and to $1$ for identical partitions (see Section~\ref{subsec:adjusted_rand_index}).}]
  {ari}{ARI}{Adjusted Rand Index}
\newacronym[description={The mean pairwise Adjusted Rand Index between the latent clusterings of independently trained surrogate policies (see Section~\ref{subsec:clustering_universality}).}]
  {cu}{CU}{\ClusteringUniversalityName}
\newacronym[description={A metric measuring how far latent clusterings differ from those of the observations, of the action distribution and of untrained networks (see Section~\ref{subsec:abstraction_score}).}]
  {as}{AS}{\AbstractionScoreName}
\newacronym[description={The Clustering Universality minus the similarity of the latent clusterings to the three trivial baselines of the Abstraction Score (see Section~\ref{subsec:adjusted_clustering_universality}).}]
  {acu}{ACU}{\AdjustedClusteringUniversalityName}

% Computing
\newacronym[description={A processor specialized in massively parallel computation, used to accelerate the training of neural networks.}]
  {gpu}{GPU}{Graphics Processing Unit}
